\scanpage{261}
\mspage{133 bis}

\uline{Prop. 9.12} : Soit $W$ un localiseur fondamental (sans plus),
i.e satisfaisant $\mathrm{Loc}(1,2,3)$. \deleted{Pour que} Les conditions
suivantes sont équivalentes

a) $W$ satisfait l'axiome de connexité $\mathrm{Loc}(8)$, i.e.
$f\in W\Longrightarrow\pi_0(f)$ bijectif.

b) Si $f,g:X\mathrel{\unsure{=|\rightrightarrows}}Y$ dans $\mathrm{Cat}$ sont telles que $\mathrm{hot}_{W}(f)=\mathrm{hot}_{W}(g)$
alors $\pi_0(f)=\pi_0(g)$.

c) \deleted{Il \unsure{y a}} existe $X,Y$ dans $\mathrm{Cat}$ tel que \inserted{$\mathrm{hot}_{W}^{X,Y}:$} $\mathrm{Hom}_{\mathrm{Cat}}(X,Y)\to\mathrm{Hom}_{W}(X,Y)$ non \unsure{constante|constant}.

c') La sous-catégorie $\mathrm{Hot}_{W}^{\unsure{*|\wedge}}$ \deleted{\unsure{de}} de $\mathrm{Hot}_{W}$, formée des
objets \deleted{qui} non vides, n'est pas équivalente à la catégorie
ponctuelle.

c'') La catégorie $\mathrm{Hot}_{W}$ n'est pas équivalente à la
catégorie $e$ ni à la catégorie $\Delta^1$.

d\deleted{\unsure{'}}) L'\unsure{ens} des classes d'isomorphie d'objets de
$\mathrm{Hot}_{W}$ est infini.

\uline{Dém} On a $a\Longrightarrow b$, puisque a) équivaut à dire
que le foncteur $\pi_0:\mathrm{Cat}\to\mathrm{Ens}$ se factorise
par $\mathrm{Hot}_{W}$. D'ailleurs le foncteur correspondant

\[
 \mathrm{Hot}_{W}\longrightarrow(\mathrm{Ens})
\]

est ess. surjectif -- il y a un inverse à isom
près, $I\unsure{-|\to}\mathrm{hot}_{W}(I)$. Cela prouve que
a) implique \unsure{d')|d)}, et on a $\unsure{d'|d}\Rightarrow\unsure{c'|d'},\,\unsure{d'|d}\Rightarrow\unsure{c''|d''}$
tautologiquement. Un argument $\pm$ tautologique
montre d'ailleurs que si $\mathrm{Hot}_{W}^{*}$ est équiv. à $e$,
alors $\mathrm{Hot}_{W}$ est équivalente à $e$ ou à $\Delta^1$
(car $\varnothing$ est objet initial de $\mathrm{Hot}_{W}$) -- donc $c''\Rightarrow c'$.
Donc on a $\pm$ tautologiquement

\diagram{\begin{tikzcd}[ampersand replacement=\&,column sep=small,row sep=small]
a \arrow[r,Rightarrow] \arrow[d,Rightarrow] \& d \arrow[r,Rightarrow] \& c'' \arrow[r,Rightarrow] \& c' \\
b \arrow[rr,Rightarrow] \arrow[rrru,Rightarrow] \arrow[rrru,draw=black,decorate,decoration={snake,amplitude=1pt,segment length=7pt},no head,"\deleted{\vphantom{\Rightarrow}}"{inner sep=0pt}] \& \& c \arrow[ur,Rightarrow] \&
\end{tikzcd}}

où les implications $b\Rightarrow c\nearrow c'$ \deleted{et celle} \inserted{sont elles}
aussi triviales. \unsure{Donc|D'où}
% archive sheet number: 260
